\par\Needspace{6\baselineskip}\medskip
\begin{center}\hypertarget{native-200-article150}{}\pdfbookmark[2]{150.}{article-150}150.\end{center}
Wir veranschaulichen dieses Verfahren nun an einigen ausführlich durchgerechneten Beispielen. Sie werden zugleich sehr deutlich zeigen, wie weit sein Anwendungsbereich reicht und wie leicht und rasch es stets zum gewünschten Ziel führt\footnote{Es ist ungenau gesprochen, wenn man ein Verfahren für \textit{mehr oder weniger exakt} erklärt als ein anderes. Eine Aufgabe kann nur durch ein Verfahren als gelöst gelten, mit dem sich zumindest jeder beliebige Grad an Genauigkeit erreichen lässt. Ein Verfahren ist daher einem anderen nur insofern überlegen, als sich \textit{derselbe} Grad an Genauigkeit mit dem einen schneller und mit geringerem Aufwand, mit dem anderen langsamer und mit größerem Aufwand erreichen lässt.}.

Unser \textit{erstes} Beispiel liefert der neue Planet Juno. Dazu wählen wir die folgenden in Greenwich angestellten Beobachtungen, die uns der berühmte \textsc{Maskelyne} mitgeteilt hat.
\SourceLeaf{201}
\begin{center}
\small\setlength{\tabcolsep}{5pt}\renewcommand{\arraystretch}{1.2}
\begin{tabular}{|r|c|c|}
\hline
Mittl. Greenwich-Zeit & Scheinb. Rektasz. & Scheinb. südl. Dekl.\\
\hline
1804 Okt. 5.\ $10^{\mathrm{h}}\,51^{\mathrm{m}}\,6^{\mathrm{s}}$ & $357^{0}\,10'\,22''{,}35$ & $6^{0}\,40'\,8''$\\
17.\quad $9\quad58\quad10$ & $355\quad43\quad45{,}30$ & $8\quad47\quad25$\\
27.\quad $9\quad16\quad41$ & $355\quad11\quad10{,}95$ & $10\quad2\quad28$\\
\hline
\end{tabular}
\end{center}

Aus den Sonnentafeln ergibt sich für dieselben Zeiten
\begin{center}
\footnotesize\setlength{\tabcolsep}{3pt}\renewcommand{\arraystretch}{1.2}
\begin{tabular}{|r|c|c|c|c|c|}
\hline\noalign{\vskip 2pt}
 & \shortstack{Sonnenlänge ab\\scheinb. Äquin.} & Nutation & \shortstack{Abstand von\\der Erde} & \shortstack{Sonnen-\\breite} & \shortstack{Scheinb. Schiefe\\der Ekliptik}\\
\hline
Okt. 5 & $192^{0}\,28'\,53''{,}71$ & $+15''{,}43$ & $0{,}998\,8839$ & $-0''{,}49$ & $23^{0}\,27'\,59''{,}48$\\
17 & $204\quad20\quad21{,}54$ & $+15{,}51$ & $0{,}995\,3968$ & $+0{,}79$ & $59{,}26$\\
27 & $214\quad16\quad52{,}21$ & $+15{,}60$ & $0{,}992\,8340$ & $-0{,}15$ & $59{,}06$\\
\hline
\end{tabular}
\end{center}

Wir setzen die Rechnung so an, als wäre die Bahn noch völlig unbekannt. Deshalb müssen wir die Parallaxe auf die Orte der Erde übertragen, da wir die Orte Junos noch nicht von ihr befreien können. Zunächst reduzieren wir daher die beobachteten Orte selbst mit der scheinbaren Schiefe vom Äquator auf die Ekliptik und erhalten die folgenden Werte.
\begin{center}
\small\setlength{\tabcolsep}{6pt}\renewcommand{\arraystretch}{1.2}
\begin{tabular}{|r|c|c|}
\hline\noalign{\vskip 2pt}
 & \shortstack{Scheinb. Länge\\Junos} & \shortstack{Scheinb. Breite\\Junos}\\
\hline
Okt. 5 & $354^{0}\,44'\,54''{,}28$ & $-4^{0}\,59'\,31''{,}59$\\
17 & $352\quad34\quad44{,}51$ & $-6\quad21\quad56{,}25$\\
27 & $351\quad34\quad51{,}57$ & $-7\quad17\quad52{,}70$\\
\hline
\end{tabular}
\end{center}

Mit dieser Rechnung verbinden wir gleich die Bestimmung der Länge und Breite des Zenits am Beobachtungsort für die drei Beobachtungen. Die Rektaszension stimmt mit der Rektaszension Junos überein (da die Beobachtungen im Meridian selbst angestellt wurden), und die Deklination ist gleich der Polhöhe $=51^{0}\,28'\,39''$.
So erhalten wir
\begin{center}
\small\setlength{\tabcolsep}{6pt}\renewcommand{\arraystretch}{1.2}
\begin{tabular}{|r|c|c|}
\hline
 & Länge des Zenits & Breite\\
\hline
Okt. 5 & $24^{0}\,29'$ & $46^{0}\,53'$\\
17 & $23\quad25$ & $47\quad24$\\
27 & $23\quad1$ & $47\quad36$\\
\hline
\end{tabular}
\end{center}
\noindent\SourceLeaf{202}Nach den in Art. 72 gegebenen Regeln bestimmen wir nun die fiktiven Orte der Erde in der Ekliptikebene selbst, an denen der Himmelskörper genauso erschienen wäre wie an den wirklichen Beobachtungsorten. Setzen wir die mittlere Sonnenparallaxe $=8''{,}6$, so ergibt sich
\begin{center}
\small\setlength{\tabcolsep}{4pt}\renewcommand{\arraystretch}{1.2}
\begin{tabular}{|r|c|c|c|}
\hline
 & Längenreduktion & Abstandsreduktion & Zeitreduktion\\
\hline
Okt. 5 & $-22''{,}39$ & $+0{,}000\,3856$ & $-0^{\mathrm{s}}{,}19$\\
17 & $-27{,}21$ & $+0{,}000\,2329$ & $-0{,}12$\\
27 & $-35{,}82$ & $+0{,}000\,2085$ & $-0{,}12$\\
\hline
\end{tabular}
\end{center}
\noindent Die Zeitreduktion ist nur angegeben, um zu zeigen, dass sie überhaupt nicht merklich ist.

Anschließend sind alle Längen des Planeten wie der Erde auf das mittlere Frühlingsäquinoktium einer bestimmten Epoche zu reduzieren. Wir wählen dafür den Beginn des Jahres 1805. Nach Abzug der Nutation ist daher noch die Präzession zu addieren, die für die drei Beobachtungen jeweils $11''{,}87$, $10''{,}23$, $8''{,}86$ beträgt. Für die erste Beobachtung ist somit $-3''{,}56$ zu addieren, für die zweite $-5''{,}28$ und für die dritte $-6''{,}74$.

Schließlich sind Junos Längen und Breiten von der Aberration der Fixsterne zu befreien. Nach den bekannten Regeln sind von den Längen jeweils $19''{,}12$, $17''{,}11$, $14''{,}82$ abzuziehen und zu den Breiten $0''{,}53$, $1''{,}18$, $1''{,}75$ zu addieren. Durch diese Addition werden die Beträge kleiner, da südliche Breiten als negativ gelten.
